\documentclass[a4paper,12pt]{article}

% =============================================================================
%  数维杯（ShuWei Cup）英文稿模板 · Summary Sheet + 六章骨架
%
%  来源：2025 数维杯秋季赛 A 题优秀论文 ASHUWEICUP2518402 的骨架拆解。
%  只借结构，不借结论：原稿的 10 处硬伤（摘要自相矛盾、引用错配、正文 0 表、
%  敏感性系数无定义、附录 51 页等）在本模板里逐条改掉，改法见 README.md。
%
%  与原稿的差异（都是刻意的）：
%    - 目录带页码（原稿只有点线，目录等于没用）
%    - 图题比正文小一号（原稿图题与正文同号）
%    - 章标题 16pt / 节 13pt（原稿 18/14，标题吃三行）
%    - 符号表两列、单位写在含义里（原稿列宽截断成 “Descriptio”，且无单位）
%    - 正文内置 4 张表骨架：参数 / 三问指标对照 / 约束核验 / 敏感性排序
%    - 附录带附件说明表（原稿 51 页代码，无文件名、无运行说明）
%    - Summary Sheet 只印一次（原稿重复印了两遍）
%
%  编译：xelatex main.tex 连跑两遍（第二遍目录页码才准）
% =============================================================================

% ======== Packages ========
\usepackage[top=1in,bottom=1in,left=1in,right=1in]{geometry}
\usepackage{amsmath,amssymb}
\usepackage{fontspec}
\usepackage{titlesec}
\usepackage{fancyhdr}
\usepackage{booktabs}
\usepackage{array}
\usepackage{longtable}
\usepackage{float}
\usepackage{caption}
\usepackage{xcolor}
\usepackage{listings}
\usepackage{graphicx}
\usepackage{lastpage}
\usepackage{hyperref}
\hypersetup{hidelinks}   % 不加这句，每个 \ref/\cite 都套红框

% ======== Fonts ========
\setmainfont{Times New Roman}
\setsansfont{Arial}
\setmonofont{Courier New}

% ======== Variables（S0 一次性填好）========
\newcommand{\ContestYear}{2026}
\newcommand{\ContestName}{ShuWei Cup}   % 换赛事改这里：ShuWei Cup / MCM/ICM ...
\newcommand{\TeamNumber}{0000}
\newcommand{\ProblemChosen}{A}
\newcommand{\PaperTitle}{[Paper Title: Model and Control Strategy for ...]}
\newcommand{\KeywordsText}{[Method]; [Object]; [Indicator]; [Strategy]}
% 摘要正文放 summary.tex：\newcommand 不是 \long，正文里一分段就报
% 「Paragraph ended before \new@command was complete」，且日志不指向摘要
\newcommand{\SummaryText}{% =============================================================================
%  Summary Sheet 正文：直接写段落，段与段之间空行分隔。
%  不要塞回 main.tex 的 \newcommand —— 宏不是 \long，一分段就报
%  「Paragraph ended before \new@command was complete」，且日志不指向摘要。
%
%  五段结构（一页装下）：
%    ① 背景：现象 → 危害 → 本类方案为什么适用（4–6 句，不写方法）
%    ②③④ 每题一段：做法一句 → 关键数值 2–3 组 → 结论一句
%    ⑤ 收尾：模型与策略的整体价值（两句）
%
%  三条铁律：
%    1. 同一指标全文一个量纲；摘要里的每个数都要能在正文找到同一个数。
%    2. 数字方向与结论同向——指标变大就解释为什么，不许写 decreases。
%    3. 抑制率 = (关 − 开)/关，必须能用摘要里的两个数复算出来。
% =============================================================================

[Summary: one paragraph of background; then one paragraph per problem --- the method in
one sentence, two or three numerical results with units, and the conclusion that follows;
then a closing paragraph. 300--500 words, no formulas, no figures, no citations.]
}
\newcommand{\TotalPages}{\pageref{LastPage}}

% ======== 正文版式 ========
\linespread{1.12}
\setlength{\parindent}{1.5em}
\setlength{\parskip}{0.2em}
\captionsetup{font=small,labelsep=quad,skip=6pt}
\setlength{\tabcolsep}{4pt}   % 默认 6pt，5 列表会顶穿版心   % 图题/表题比正文小一号，且无冒号

% ======== 标题层级（原稿 18/14 太大，收到 16/13）========
\renewcommand{\thesection}{\arabic{section}.}
\renewcommand{\thesubsection}{\arabic{section}.\arabic{subsection}}
\titlespacing*{\section}{0pt}{1.4em}{0.7em}
\titlespacing*{\subsection}{0pt}{1.0em}{0.45em}
\titleformat{\section}[block]{\fontsize{16pt}{19pt}\bfseries}{\thesection}{0.6em}{}
\titleformat{\subsection}[block]{\fontsize{13pt}{16pt}\bfseries}{\thesubsection}{0.5em}{}

% ======== 页眉：Team #0000 ／ Page N of M ========
% 原稿印成 “Page3of 23”，缺空格；这里用 \thepage\ of 让 TeX 自己管间距
\setlength{\headheight}{16pt}
\pagestyle{fancy}
\fancyhf{}
\fancyhead[L]{\sffamily Team \#\TeamNumber}
\fancyhead[R]{\sffamily Page \thepage\ of \TotalPages}
\renewcommand{\headrulewidth}{0.4pt}

% ======== 三线表 / 图 / 代码 ========
% #1 caption  #2 label  #3 colspec  #4 header  #5 body
\newcommand{\threelinetable}[5]{%
  \begin{table}[htbp]
    \centering
    \caption{#1}
    \label{#2}
    \begin{tabular}{#3}
      \toprule
      #4 \\
      \midrule
      #5 \\
      \bottomrule
    \end{tabular}
  \end{table}%
}
% \fig{宽度比例}{图片路径}{图题}{label}
\newcommand{\fig}[4]{%
  \begin{figure}[htbp]
    \centering
    \includegraphics[width=#1\textwidth]{\detokenize{#2}}
    \caption{#3}
    \label{#4}
  \end{figure}%
}
\lstset{
  frame=topbottom,
  framerule=0.55pt,
  framesep=0pt,
  basicstyle=\ttfamily\small,
  backgroundcolor=\color[gray]{0.97},
  numbers=none,
  showstringspaces=false,
  breaklines=true,
}

% ======== Document ========
\begin{document}

% ---------------------------------------------------------------- Summary Sheet
% 一页装下全部结论：三栏表头 + 标题 + 5 段摘要 + 关键词。
% 原稿在这里犯的错：问题一 未启动 95.4 → 启动 118.6（更差）却写 "significantly
% decreases"。摘要的每个数字都必须与结论同向，写完逐段对一遍。
\thispagestyle{empty}
\vspace*{-1.2em}
\noindent
\begin{minipage}[t]{0.29\textwidth}\centering
  \textbf{Question Chosen}\\[0.5em]
  {\fontsize{20pt}{22pt}\selectfont\bfseries\textcolor{red!75!black}{\ProblemChosen}}
\end{minipage}%
\begin{minipage}[t]{0.40\textwidth}\centering
  \textbf{\ContestYear\ \ContestName}\\[0.25em]
  \textbf{Summary Sheet}
\end{minipage}%
\begin{minipage}[t]{0.29\textwidth}\centering
  \textbf{Team Control Number}\\[0.5em]
  {\fontsize{16pt}{18pt}\selectfont\bfseries\textcolor{red!75!black}{\TeamNumber}}
\end{minipage}
\vspace{0.9em}

\noindent\rule{\textwidth}{0.8pt}
\vspace{1.6em}

\begin{center}
  {\fontsize{16pt}{19pt}\selectfont\bfseries \PaperTitle}
\end{center}
\vspace{1.2em}

\SummaryText

\vspace{0.8em}
\noindent\textbf{Keywords:} \KeywordsText
\newpage

% ---------------------------------------------------------------- Contents
% 原稿目录只有点线、没有页码（等于没有目录）。这里交给 \tableofcontents，
% 页码随重建自动正确——不要手写目录。
\pagenumbering{roman}
\setcounter{page}{1}
\begin{center}{\fontsize{16pt}{19pt}\selectfont\bfseries Contents}\end{center}
\vspace{1.2em}
\renewcommand{\contentsname}{}
\tableofcontents
\newpage

% ---------------------------------------------------------------- Body
\pagenumbering{arabic}
\setcounter{page}{1}
\pagestyle{fancy}

\section{Introduction}

\subsection{Background}

%\noindent 写这一段只做三件事：现象是什么、危害在哪、为什么这一类方案适用。
% 不要在这里展开模型，也不要提前给结果——摘要已经给过了。
% 长度：4–6 句，一段到底，不分小点。
Large vehicles experience lateral vibration during operation ...
(symptom) ... It accelerates fatigue of the mounting structure and degrades
(positioning / comfort) ... Among the active schemes, centrifugal actuators
are suitable here because (why this class fits this problem) ...

\subsection{Work}

% 逐问复述"这道题要交什么"——一小节一问。
% 好处：评委读完 1.2 就知道三问各要交哪些曲线、哪些指标，
% 后面读到结果时不用回头找题目要求。
\subsubsection{Problem 1}
Under the operating conditions of Scenario 1 ... establish the lateral model,
report the displacement and acceleration curves for $0$--$10$\,s with the actuator
inactive and active, and give the vibration index for both states.

\subsubsection{Problem 2}
Under the multi-source disturbance of Scenario 2 ... design the angular displacement
schedule that minimises the index, report both states' curves and the schedule itself,
and verify the angular velocity / acceleration limits.

\subsubsection{Problem 3}
With the disturbance unmeasurable and only the lateral acceleration observable ...
design a schedule that minimises the index from the measurable signals alone, and
report the index for both states.

% ┈┈┈┈┈┈┈┈┈┈┈┈┈┈┈┈┈┈┈┈┈┈┈┈┈┈┈┈┈┈┈┈┈┈┈
% 【Section writing prompt】
% ════════════════════════════════════════════
% Position: the first section. 1.1 background is one paragraph, 1.2 repeats the tasks.
% 1.2 is not a summary of the paper: it restates what each question asks for. Write it in
%           the problem's own vocabulary, keep one subsubsection per question.
% Do not:   put method names in 1.1 unless the method is the subject; put results in 1.2.
% Citation: 1–2 references at most in 1.1 (the context claim), none in 1.2.
% ════════════════════════════════════════════

\section{Question Analysis}

\subsection{Data Analysis}
The raw table carries columns beyond the model's needs, so the target variable is first
located and matched to the physical quantity the model expects. [Then list each
operation and its reason: fill / truncate / interpolate.]

\subsection{Analysis of Problem 1}
[Known, unknown, framework, output — four beats in one paragraph.]

\subsection{Analysis of Problem 2}
[Same four beats; name the key difficulty in one sentence.]

\subsection{Analysis of Problem 3}
[Same four beats; state what information is missing.]

% 【Section writing prompt】
% 2.1 每个处理都要给"为什么"，不是报流水账：只写"缺失值填 0"没人信，写清
% "填 0 表示无输入，是保守策略，避免注入不合理扰动"才有说服力。
% 2.2–2.4 一段一问：已知 → 未知 → 框架 → 输出，再加一句"关键难点在哪"。
% 能从句面直接抄回题面的句子，删掉。

\section{Symbols and Assumptions}

\subsection{Symbols}

The notation used throughout the paper is fixed in Table~\ref{tab:symbols}: rows follow
the order of first appearance, one symbol means exactly one quantity, and the unit is
written in parentheses inside the meaning cell.

\begin{longtable}{@{}>{\centering\arraybackslash}p{0.16\textwidth}>{\raggedright\arraybackslash}p{0.80\textwidth}@{}}
  \caption{Symbols.}
  \label{tab:symbols} \\
  \toprule
  Symbol & Meaning \\
  \midrule
  \endfirsthead
  \caption{Symbols (continued).} \\
  \toprule
  Symbol & Meaning \\
  \midrule
  \endhead
  \bottomrule
  \endfoot
  $x_1$ & [quantity] (unit)\\
  $x_2$ & [quantity] (unit)\\
  $y(t)$ & [quantity] (unit)\\
  $\theta$ & [quantity] (rad)\\
  $\Delta t$ & sampling step (s)\\
  $I$ & [objective / index] (unit)\\
\end{longtable}

\subsection{Assumptions}
\begin{enumerate}
  \item [Simplification assumption shared by all questions.]
  \item [Parameter / modelling assumption shared by all questions.]
  \item \textbf{(Problem 2)} [What the scenario fixes: knowns, decision variable.]
  \item \textbf{(Problem 3)} [What is unmeasurable, what replaces it.]
\end{enumerate}

% 【Section writing prompt】
% 位置：Question Analysis 之后、Model 之前——读者必须先见到每个符号，才遇到用到它的公式。
% 符号表：两列，单位写在含义里；一个符号一个含义；循环下标不进表；数值放参数表。
% 假设：场景相关的那几条要标出服务哪一问（如上例 3、4），套话式假设（"本文假设…合理"）全部删。

\section{Model}

\subsection{Model of Problem 1}

\paragraph{Force analysis.} The body is acted on by three forces: the lateral disturbance
$F_d(t)$, the spring--damper constraint force, and the actuator resultant force. Taking
the lateral direction as positive,
\begin{equation}
  F_h(t) = -k\,y(t) - c\,\dot y(t),
  \label{eq:fh}
\end{equation}
with $k$ and $c$ listed in Table~\ref{tab:params}.

\paragraph{Actuator force.} A single eccentric block of radius $r$ rotating at
$\omega(t)$ exerts $m r \omega^2$ along its radius. The actuator carries two blocks
rotating in opposite directions at equal speed, so the lateral components add and the
in-plane components cancel; for two actuators the resultant is
\begin{equation}
  F_a(t) = 2mr\left[\omega_i^2\cos\theta_i + \alpha_i\sin\theta_i\right] \quad\text{(per actuator pair)}.
  \label{eq:fa}
\end{equation}

\paragraph{Equation of motion.} Substituting \eqref{eq:fh} and \eqref{eq:fa} into Newton's
second law for the lateral translation gives
\begin{equation}
  M\ddot y(t) + c\dot y(t) + k y(t) = F_d(t) + F_a(t),
  \label{eq:eom}
\end{equation}
which is integrated with a fourth-order Runge--Kutta scheme at $\Delta t$, with the
angular velocity and acceleration obtained from $\theta_i(t)$ by finite differences.

\paragraph{Vibration index.} With $T = 10$\,s,
\begin{equation}
  I_h = \frac{1}{T}\int_0^{T} a_y^2(t)\,\mathrm{d}t,
  \label{eq:index}
\end{equation}
evaluated by trapezoidal integration.

% 参数表：一次写清，后面正文不再重复数字。
\threelinetable{Model parameters.}{tab:params}{@{}p{0.34\textwidth}lrr@{}}
  {Quantity & Symbol & Value & Unit}
  {Body mass & $M$ & 2000 & kg \\
   Eccentric mass & $m$ & 100 & kg \\
   Eccentric radius & $r$ & 0.2 & m \\
   Lateral stiffness & $k$ & $7.2252\times10^{6}$ & N/m \\
   Lateral damping & $c$ & 3600 & N$\cdot$s/m \\
   Angular velocity limit & $\omega_{\max}$ & 100 & rad/s \\
   Angular acceleration limit & $\alpha_{\max}$ & 5000 & rad/s$^2$}

\subsection{Result of Problem 1}

% 结果先核约束、再报指标——这是原稿做得最好的一个动作，务必保留。
% 只报"降了 80 %"却不证明作动器没超限，评委会怀疑方案不可执行。
With the actuator inactive the index is $I_h^{\text{off}} = [\text{value}]$ [unit]; with the
actuator driven by the prescribed schedule it is $I_h^{\text{on}} = [\text{value}]$ [unit].
Table~\ref{tab:limits} confirms the schedule stays inside the actuator's envelope, so the
comparison is meaningful.

\begin{table}[htbp]
  \centering
  \caption{Constraint check and index summary for the three problems.}
  \label{tab:limits}
  \small
  \begin{tabular}{@{}lrrrr@{}}
    \toprule
    Case & $I_h$ off & $I_h$ on & Peak $\omega$ (rad/s) & Peak $\alpha$ (rad/s$^2$) \\
    \midrule
    Problem 1 & [ ] & [ ] & [ ] & [ ] \\
    Problem 2 & [ ] & [ ] & [ ] & [ ] \\
    Problem 3 & [ ] & [ ] & [ ] & [ ] \\
    \midrule
    Limit     & --  & --  & 100 & 5000 \\
    \bottomrule
  \end{tabular}
\end{table}

\subsection{Model of Problem 2}

% 与问题一的差别只写"差在哪"，不重抄力分析。原稿这一点做对了。
The force analysis is unchanged; the difference is that $\theta_i(t)$ is now the decision
variable. Over a horizon of $N_p$ steps, at each step solve
\begin{equation}
  \min_{\{\theta_i\}} \; \hat I_h
  \quad\text{s.t.}\quad
  M\ddot y + c\dot y + k y = F_d + F_a(\theta,\omega,\alpha),\;\;
  |\omega|\le \omega_{\max},\;\; |\alpha|\le \alpha_{\max},
\end{equation}
apply the first step, and repeat.

\subsection{Result of Problem 2}

Report the index for both states, the suppression rate, the schedule itself (Figure), and
the peak $\omega$ / $\alpha$ against their limits.

\subsection{Model of Problem 3}

% 问题三的重点是"少了什么信息、用什么补回来"，不是再推一遍动力学。
With $F_d$ unmeasurable, only $a_y$ is observed. An extended state observer reconstructs
$y$, $\dot y$ and the lumped disturbance $z$:
\begin{equation}
  \dot{\hat y} = \hat v + \beta_0 (y - \hat y),\quad
  \dot{\hat v} = a_y + \beta_1 (y - \hat y),\quad
  \dot{\hat z} = \beta_2 (y - \hat y),
\end{equation}
and a PID loop on the reconstructed state drives the angular displacement:
\begin{equation}
  \theta(t) = K_p e + K_i\!\int e\,\mathrm{d}t + K_d \dot e .
\end{equation}

\subsection{Result of Problem 3}

Report the index for both states, the suppression rate, the disturbance estimation error
after convergence, and the gain sensitivity (Section 5.3).

% ┈┈┈┈┈┈┈┈┈┈┈┈┈┈┈┈┈┈┈┈┈┈┈┈┈┈┈┈┈┈┈┈┈┈┈
% 【Section writing prompt】
% ════════════════════════════════════════════
% Position: the body of the paper; it should take about 60 % of the page budget.
% Five beats per question, identical across questions: force analysis → equations
%           (numbered) → numerical method → results → index.
% Results:  check the actuator limits BEFORE quoting the index. A suppression rate with
%           no feasibility check invites the reader to doubt the scheme.
% Tables:   at least three — parameters, index summary, constraint check. A paper whose
%           results exist only as curves cannot be audited.
% Do not:   re-derive what Problem 1 already derived. State the difference and move on.
% Units:    one unit per quantity for the whole paper; the index unit must match the
%           dimension of the defining equation, not a plausible-looking one.
% ════════════════════════════════════════════

\section{Sensitivity Analysis}

% 系数必须先定义再引用，否则读者无法解释——这是这类章节最常见的漏洞。
Define the relative sensitivity of the index to parameter $p$ as
\begin{equation}
  S_p = \frac{\Delta I / I_0}{\Delta p / p_0},
  \label{eq:sensitivity}
\end{equation}
where $I_0$ is the baseline index at the nominal set and every perturbation is
$\pm 10\,\%$ about it.

\subsection{Sensitivity of Problem 1}
[Rank $|S_p|$; explain any peak by the mechanism, not by the number.]

\subsection{Sensitivity of Problem 2}
[Same procedure on the schedule's parameters.]

\subsection{Sensitivity of Problem 3}
[Observer gains and controller gains; a coefficient that rounds to zero is a finding.]

% 【Section writing prompt】
% 三件事：一个公式（定义敏感性系数）、一个固定摄动幅度（默认 ±10 %，全文统一）、一张排序表。
% 口径不许混：一问答"平均敏感系数"、另一问答"百分比变化"，两套数放一起就没法比较。
% 全为 0 的参数要下结论（"该参数不敏感"），不要罗列五个零。
% 图：三小图并排，纵轴指数、横轴 Parameter Variation (%)、虚线标 Baseline——基线让相对变化可读。

\section{Strengths and Weaknesses}

\subsection{Strengths}

% 每条挂具体机理，不写"模型简单""精度较高"这类空话。
\begin{enumerate}
  \item The force decomposition is explicit: disturbance, spring--damper constraint and
        actuator resultant force are derived separately, so each term can be checked
        against a physical measurement.
  \item The three problems share one equation of motion; only the decision variable
        changes. A reader who understands Problem 1 can follow Problem 3 without new
        mechanics.
  \item Every reported index is accompanied by a feasibility check on the angular
        velocity and acceleration limits, so the schedule is executable rather than
        merely optimal.
  \item The data handling is justified operation by operation (fill, truncate,
        interpolate), and the interpolation error is bounded by the sampling step.
  \item The sensitivity analysis is run per problem, and the stiffness peak is explained
        by the natural-frequency crossing rather than reported as a number.
\end{enumerate}

\subsection{Weaknesses}

% 不足要写"会因为什么原因、影响哪个结论"，最好带一条改进方向。
\begin{enumerate}
  \item Parameters are treated as constant; load, wear and temperature make $M$, $k$,
        $c$ time-varying, so the robustness claim holds only near the nominal set.
  \item Friction in the eccentric-block bearings and losses in the bevel-gear train are
        ignored, which biases the derived angular acceleration and hence $F_a$.
  \item The observer is tuned on the nominal model; gain scheduling for a strongly
        time-varying disturbance is left open.
  \item Only a single disturbance realisation per scenario is reported; a Monte Carlo
        sweep over disturbance phase would give the suppression rate a confidence band.
\end{enumerate}

% ┈┈┈┈┈┈┈┈┈┈┈┈┈┈┈┈┈┈┈┈┈┈┈┈┈┈┈┈┈┈┈┈┈┈┈
% 【Section writing prompt】
% ════════════════════════════════════════════
% Counts:   at least 4 strengths and 4 weaknesses. Three-and-two (the usual) reads as
%           filler.
% Every item names a mechanism and its consequence: "friction in the bearings biases the
%           angular acceleration, hence F_a" beats "the model is simplified".
% Weakness items should hint at the fix — that is what turns a weakness list into a
%           research direction.
% ════════════════════════════════════════════


% ---------------------------------------------------------------- References
\clearpage
% =============================================================================
%  References
%
%  规则（check_paper.py 会机器核）：
%   1. 正文用 \cite{key}，key 与这里一一对应；没被引用的条目等于装饰品。
%   2. 条目必须能查到：作者、题名、出处、年、卷期页。写不出完整出版信息就不要写。
%   3. 引的是"这个方法本身"。原稿的致命错误就在这里：为 MPC 引了
%      "magnesium phosphate cement (MPC)" 与 "MPC-mediated lactate production"
%      —— 两篇都不是 Model Predictive Control，评审一眼可见。
%   4. 8–15 条。少于 8 条说明没读文献，多于 15 条说明没筛。
% =============================================================================

\addcontentsline{toc}{section}{References}
\begin{thebibliography}{9}

\bibitem{ref1}
[Author A, Author B.] [Title of the paper on the modelling method used in Section 4][J].
[Journal], [year], [volume]([issue]): [pages].

\bibitem{ref2}
[Author C.] [Title on the control formulation][J]. [Journal], [year], [volume]([issue]): [pages].

\bibitem{ref3}
[Author D, Author E.] [Title on the observer / estimation technique][J]. [Journal], [year],
[volume]([issue]): [pages].

\bibitem{ref4}
[Author F.] [Title on the physical mechanism of the actuator][J]. [Journal], [year],
[volume]([issue]): [pages].

\bibitem{ref5}
[Author G, Author H.] [Title on vibration indices / evaluation of suppression][J].
[Journal], [year], [volume]([issue]): [pages].

\bibitem{ref6}
[Author I.] [Title on sensitivity analysis of dynamic systems][J]. [Journal], [year],
[volume]([issue]): [pages].

\bibitem{ref7}
[Author J, Author K.] [Title on numerical integration / solver choice][J]. [Journal],
[year], [volume]([issue]): [pages].

\bibitem{ref8}
[Author L.] [Title on an application of the same actuator class][J]. [Journal], [year],
[volume]([issue]): [pages].

\end{thebibliography}


% ---------------------------------------------------------------- Appendix
\clearpage
\section*{Appendix}
\addcontentsline{toc}{section}{Appendix}
% =============================================================================
%  Appendix — core code
%  顺序：先附件说明表，再贴代码。附表说明"哪个文件对应哪个问题、产出什么"，
%  审查者找不到映射，就不会去读代码。
%  篇幅：附录压在正文的 1/3 以内；放不下就把完整代码单独提交，正文只留核心循环。
% =============================================================================

\subsection*{A.1 File list}

\begin{table}[htbp]
  \centering
  \caption{Submitted code files.}
  \label{tab:codefiles}
  \footnotesize
  \begin{tabular}{@{}p{0.17\textwidth}p{0.15\textwidth}p{0.40\textwidth}p{0.16\textwidth}@{}}
    \toprule
    File & Problem & Purpose & Output \\
    \midrule
    \texttt{utils.py}    & all       & parameters, index evaluation, plotting & -- \\
    \texttt{problem1.py} & Problem 1 & [what it does]                         & \texttt{[file]} \\
    \texttt{problem2.py} & Problem 2 & [what it does]                         & \texttt{[file]} \\
    \bottomrule
  \end{tabular}
\end{table}

\subsection*{A.2 Core code}

\begin{lstlisting}[language=Python]
# problem1.py -- forward simulation
# run: python problem1.py   -> writes [output file]
import numpy as np
from utils import params, vibration_index

P = params()
# [core loop: one step per sample, accumulate the index]
\end{lstlisting}

% 【Appendix writing prompt】
% 每个 listing 第一行写死 "run: ... -> output"，读者不用猜怎么跑。
% 一题一份可运行脚本；草稿版、重复 import 的副本不要贴。


\end{document}
